\section{Discussion}
\label{sec:discussion}

\subsection{Key Implications}

\textbf{Decentralization is Achievable}: Holochain DHT integration demonstrates that Byzantine-robust federated learning can be fully decentralized, achieving 10,127 TPS with 89ms median latency and zero transaction costs while eliminating centralized trust assumptions. This is our primary contribution.

\textbf{Detector Inversion is Real}: The direct A/B comparison provides definitive empirical evidence that peer-comparison methods experience catastrophic failure (100\% FPR) at 35\% BFT with heterogeneous data. This validates theoretical concerns and demonstrates the necessity of external quality signals.

\textbf{Server-Side Validation Works at Low-to-Moderate BFT}: Both PoGQ and FLTrust achieve perfect detection (100\% TPR, 0\% FPR) at 20--30\% Byzantine ratios, confirming that server-side validation is effective for realistic threat models.

\subsection{Limitations of Loss-Based Quality Metrics (PoGQ)}

While PoGQ achieves 100\% detection at 20--35\% BFT on MNIST and 20--30\% on FEMNIST, our evaluation reveals performance degradation at higher Byzantine ratios compared to direction-based methods like FLTrust.

\textbf{FEMNIST High-BFT Performance}: Across three random seeds (42, 123, 456), PoGQ exhibits inconsistent performance above 30\% BFT. At 35\% BFT, detection rates range from 71--100\% (mean: 90.3\%, $\sigma$: 16.8\%), and at 45--50\% BFT, PoGQ fails completely (0\% TPR) on most seeds. In contrast, FLTrust maintains 100\% TPR across all tested ratios including 50\% BFT.

\textbf{CIFAR-10 Failure}: Grid search across 193 hyperparameter configurations on CIFAR-10 with ResNet18 revealed systematic failure---all configurations achieved 0\% detection rate with AUROC < 0.5 (worse than random). We hypothesize this stems from increased gradient diversity in high-dimensional parameter spaces (ResNet18: 11M parameters vs. SimpleCNN: 50K), where honest gradients from heterogeneous clients produce highly varied loss improvements, causing distribution overlap with Byzantine quality scores.

\textbf{Root Cause Hypothesis}: Loss-based metrics measure gradient \emph{utility} on validation data. In high-BFT regimes with heterogeneous client data, Byzantine nodes can produce gradients that improve loss on narrow validation subsets, creating overlap with honest quality score distributions. Direction-based metrics (like FLTrust's cosine similarity) appear more robust to this phenomenon, as gradient direction is less sensitive to data heterogeneity than absolute loss improvement.

\textbf{Practical Implications}: For deployments requiring robustness at 40--50\% BFT, we recommend FLTrust over PoGQ. PoGQ remains competitive at 20--30\% BFT and may be preferable when direction-based trust scoring is insufficient (e.g., optimization-based poisoning attacks that preserve gradient direction).

\subsection{General Limitations and Assumptions}

\textbf{Server Trust}: Mode 1 assumes the server possesses a clean validation set. This is reasonable in many federated learning scenarios (e.g., healthcare institutions, enterprises) but limits applicability to fully trustless environments. \textit{Mitigation strategies}: (1) \textbf{Federated validation} via secret sharing~\cite{bonawitz2017practical} distributes validation across multiple nodes, (2) \textbf{Zero-knowledge proofs} (zk-SNARKs~\cite{bensasson2014succinct}) enable cryptographic verification without revealing validation data, (3) \textbf{Multi-party computation} allows collaborative quality assessment without central oracle.

\textbf{Computational Overhead}: PoGQ requires $O(N)$ forward passes through the model per round (2--7\% overhead). For very large networks ($N > 1000$), sampling strategies may be necessary. \textit{Optimization}: Holochain's DHT naturally distributes validation across $\sqrt{N}$ validators per entry, reducing per-node overhead. For large gradient storage, IPFS content-addressing with Filecoin incentivized persistence can offload DHT storage costs while maintaining tamper-evidence via cryptographic hashes.

\textbf{Data Heterogeneity}: Our experiments use Dirichlet($\alpha$=0.1) label skew. More extreme heterogeneity (e.g., $\alpha$=0.01) may increase quality score overlap, potentially raising FPR above 10\%.

\textbf{Attack Models}: We evaluate sign flip, scaling, noise, and sleeper agents. More sophisticated attacks (e.g., gradient clipping evasion, optimization-based poisoning) require further study.

\subsection{Practical Deployment Considerations}

\textbf{Validation Set Size}: Our 1000-sample validation set provides stable quality scores. Smaller sets (100--500 samples) may increase variance.

\textbf{Threshold Tuning}: The adaptive threshold algorithm requires no manual tuning but assumes bimodal quality score distribution (Byzantine vs. honest clusters). Pathological distributions may require domain-specific adjustments.

\textbf{Network Dynamics}: Node churn, network partitions, and asynchronous communication require additional mechanisms beyond our current scope.

\subsection{Future Work}

\textbf{Improved Loss-Based Metrics}: Address CIFAR-10 failure through: (1) \textbf{Gradient decomposition}---assess quality per layer/module rather than aggregate loss improvement, (2) \textbf{Multi-metric fusion}---combine loss improvement with direction similarity and parameter magnitude, (3) \textbf{Adaptive validation}---dynamically select validation samples most sensitive to Byzantine attacks.

\textbf{Zero-Knowledge Proofs}: Integrate zk-SNARKs for cryptographic gradient quality proofs, eliminating server trust assumption.

\textbf{Federated Validation}: Distribute validation set across multiple honest nodes using secret sharing, reducing single-point-of-trust.

\textbf{Adaptive Sampling}: For large networks ($N > 1000$), sample subset of gradients for validation using reputation-weighted strategies.

\textbf{Advanced Attacks}: Evaluate against optimization-based poisoning~\cite{bagdasaryan2020backdoor} and gradient clipping evasion.

\textbf{Additional Datasets}: Extend to (1) \textbf{Tabular data} (healthcare, finance) with feature heterogeneity, (2) \textbf{Non-image domains} (NLP sentiment analysis, time series forecasting) to validate decentralized architecture generality.

\subsection{Security vs Performance Tradeoffs}

We adopt RISC Zero zkVM (v3.0) as our primary proof backend, providing 127-bit conjectured security (S128 profile) with 30--60 second proving time on consumer hardware. While domain-specific AIRs like Winterfell theoretically offer 7--10$\times$ performance improvements, they require static trace length determination at compile-time. PoGQ's variable-length gradient vectors (dynamic based on model architecture) fundamentally conflict with this constraint. We explored selector columns and padding strategies but could not resolve this without sacrificing soundness or generality.

RISC Zero's general-purpose RISC-V execution model eliminates this limitation, allowing identical PoGQ decision logic to handle arbitrary gradient dimensions without AIR recompilation. The 30--60 second proving cost is acceptable for federated learning audit scenarios where: (1) verification is instantaneous (<100ms), (2) proofs are generated asynchronously per-round (not blocking training), (3) cryptographic auditability justifies computational overhead. For higher security requirements, S192 profile provides 191-bit security at 2$\times$ proving cost. This architecture prioritizes correctness and flexibility over raw performance, aligning with our threat model where Byzantine detection accuracy matters more than proving throughput.

\subsection{Broader Impact}

Zero-TrustML demonstrates that decentralized Byzantine-robust federated learning is practical, enabling new deployment scenarios: IoT networks with compromised devices, open federated learning platforms, cross-institutional collaboration with partial trust. By eliminating centralized aggregation, we expand the applicability of privacy-preserving collaborative ML to scenarios where centralized trust is unacceptable. Our comparative evaluation provides guidance for practitioners selecting between loss-based (PoGQ) and direction-based (FLTrust) validation methods based on expected threat models.
